
\documentclass[12pt]{ltjarticle}
\usepackage[dvips, twoside, bindingoffset=0.4cm, paper=a4paper, hmargin=1.5cm, top=3cm, bottom=1.5cm]{geometry}
\usepackage[psamsfonts]{amssymb}
\usepackage[dvips]{graphicx}

\setlength{\textheight}{25cm}
\setlength{\textwidth}{17.0cm}
\setlength{\evensidemargin}{-0.5cm}
\setlength{\oddsidemargin}{-0.5cm}

\begin{document}
%\thispagestyle{empty}
\pagestyle{empty}

\vspace*{1cm}

\begin{center}

\begin{Large}
{\bf Proceedings of the 22nd International Conference of}
\end{Large} 
\\
\begin{huge}
{\bf Logic and Engineering of \\
% \\
Natural Language Semantics 22} 
\end{huge}
\\\ \\
\begin{huge}
{\bf (LENLS22)}
\end{huge}

\vspace*{1cm}

\begin{Large}
{\em hosted by \\\ \\ The Association for Logic, Language and Information (FoLLI)}
\end{Large}

\begin{large}
\vspace*{1cm}
Workshop Chair\\ $\;$\\
Osamu Sawada (Kobe University) \\
\ \\
Workshop Co-chairs\\ $\;$\\
Daisuke Bekki (Ochanomizu University) \\
Koji Mineshima (Keio University) \\

\end{large}

\vspace*{3cm}

%\begin{center}
%\ifpdf{
%\centerline{
%\raisebox{1ex}{\scalebox{.3}{\includegraphics{jsai_logo.eps}}}
%\hspace*{30pt}
%\scalebox{.3}{\includegraphics{logo05.eps}}
%}
%}
%\ifdvi{
%\centerline{
%\includegraphics[width=3.5cm,height=2cm]{jsai_logo.eps}
%}
%}
%\end{center}

\vspace*{1cm}

\begin{large}
Kobe University \\
November 25 (Wed) - 28 (Sat), 2026
\end{large}
\vfill

\end{center}

\newpage

\vspace*{20cm}
\vfill
\begin{large}

\end{large}

\newpage
\pagestyle{plain}
\pagenumbering{roman}

\newpage
\section*{Preface}
%%%%%%%%%%%%%%%%%%
This online proceedings volume contains selected and invited papers on topics in the inter-disciplinary fields of linguistics, logic, computation, and philosophy, including the following:
\begin{itemize}
\item[$\maltese$] Formal syntax, semantics and pragmatics of natural language
\item[$\maltese$] Model-theoretic and/or proof-theoretic semantics of natural language
\item[$\maltese$] Computational approaches to semantics and pragmatics
\item[$\maltese$] Nonclassical logic and its relation to natural language 
  (e.g. substructural, fuzzy, categorical, and topological logic)
\item[$\maltese$] Formal philosophy of language
\item[$\maltese$] Scientific methodology and/or experimental design in linguistics
\end{itemize}

\subsection*{Sponsors}
LENLS is being organized by an alliance of the following sponsors.
\begin{itemize}
\item The Association for Logic, Language and Information (FoLLI)
\item "The Nexus of Human-LLM Interaction: A Multidisciplinary Linguistic Approach to Building a Collaborative Feedback Loop" project (Grant Number JPMJCR2565), funded by JST CREST Programs "Creation of an Interdisciplinary System Foundation for Realizing a Society of Human-AI Coexistence and Collaboration"
\item Graduate School of Humanities, Kobe University
\item JSPS KAKENHI Grant Number JP23H03452, Japan.
\end{itemize}

\subsection*{Workshop Organizers/Program Committee}
\begin{flushleft}
\begin{tabular}{l}
Daisuke Bekki (Ochanomizu University) \\
Alastair Butler (Hirosaki University) \\
Patrick D. Elliott (Heinrich-Heine University of Düsseldorf) \\
Naoya Fujikawa (University of Tokyo) \\
Hayate Funakura (Keio University) \\
Robert Henderson (University of Arizona) \\
Yu Izumi (Nanzan University) \\
Magdalena Kaufmann (University of Connecticut) \\
Elin McCready (ICREA/Autonomous University of Barcelona) \\
Koji Mineshima (Keio University) \\
Teruyuki Mizuno (Ochanomizu University) \\
Kenta Mizutani (Aichi Prefectural University) \\
Yoshiki Mori (University of Tokyo) \\
Yang Muyi (Osaka University) \\
Shinya Okano (University of Tokyo) \\
David Y. Oshima (Nagoya University) \\
Lukas Rieser (Tokyo University of Agriculture and Technology) \\
Osamu Sawada (Kobe University) \\
Eri Tanaka (Osaka University) \\
Wataru Uegaki (University of Edinburgh) \\
Grégoire Winterstein (Université du Québec à Montréal) \\
Shunsuke Yatabe (Ochanomizu University) \\
Kei Yoshimoto (Tohoku University) \\

\end{tabular}
\end{flushleft}
\newpage
  
\section*{Program and Table of Contents}

\newcommand{\slot}[2]{\noindent \underline{#1 \  #2} \\}
\newcommand{\talk}[3]{
  \noindent #2 \\ 
  \indent\indent \textit{#1} \dotfill #3 
  \smallskip \\
  }
\newcommand{\talkk}[3]{
  \noindent #2 \\ 
  \indent\indent \textit{#1}
  \smallskip \\
  }


\noindent\textbf{\large 
Pre-conference Event: November 25 (Wed)
}\\



\slot{13:30-18:00}{JST CREST-PRESTO Joint Symposium (separate registration)}




\noindent\textbf{\large 
1st Day: November 26 (Thu)
}\\



\slot{9:50-10:00}{Opening remark}



\slot{10:00-11:30}{Anaphora and dynamic semantics}

  
\talk{Patrick Elliott and Alexandros Kalomoiros}
     {Complement subordination}
     {pp.1--6}
  
\talk{Oleg Kiselyov}
     {From Dependent Anaphora to Paycheck Pronouns in Static Orthodox First-Order Logic}
     {pp.7--12}
  
\talk{Hajime Mori and Soichiro Kinoshita}
     {A CDRT Analysis of the Indefinite Interpretation of Russian Pronouns as Predicate Anaphora}
     {pp.13--17}



\slot{13:30-15:00}{Social meaning and perspective}

  
\talk{Elin McCready}
     {Formalizing Perspectives}
     {pp.18--22}
  
\talk{Elin McCready and Grégoire Winterstein}
     {Just Asking Questions?}
     {pp.23--27}
  
\talk{Samuel Laperle}
     {French gym bros are bad roommates: an attempt to characterize the semantics of personae}
     {pp.28--32}



\slot{15:30-17:00}{Logic and proof theory}

  
\talk{Hinari Daido and Daisuke Bekki}
     {Disjunctive Anaphora in Dependent Type Semantics and its Automated Theorem Prover}
     {pp.33--39}
  
\talk{Timothée Bernard and Philippe de Groote}
     {Two ways to be a singer: eventive and dispositional readings in perspective semantics}
     {pp.40--43}
  
\talk{Kohei Kishida}
     {Aristotle's Substructural Sequent Calculus (Extended Abstract)}
     {pp.44--49}



\slot{17:30-18:30}{Invited Talk}

  
\talk{Thomas Grano}
     {The fine-grained modal semantics of English for-to infinitives}
     {p.50}




\noindent\textbf{\large 
2nd Day: November 27 (Fri)
}\\



\slot{10:00-11:30}{Modality and pragmatics in Japanese and beyond}

  
\talk{Taro Hiraoka}
     {Monotonic and non-monotonic confirmation-seeking: A scoreboard analysis of daroo? and daroo-ne? in Japanese}
     {pp.51--56}
  
\talk{Yu Izumi}
     {A Minimal Participant-Structure Semantics for Japanese Affixal Designation Terms}
     {pp.57--62}
  
\talk{Christopher Davis, Tim Jantarungsee and Yasutada Sudo}
     {How to Build a Modal: A Semantic Typology of Conditional Evaluative Constructions}
     {pp.63--67}



\slot{13:30-15:00}{Experimental and computational approaches}

  
\talk{Yixuan Yan}
     {Dissociating Disjunctive Meaning from Clause-type Restrictions: Evidence from Child Mandarin}
     {pp.68--72}
  
\talk{Kotomi Fujita, Daniel Plesniak, Haruka Fukushima and Daisuke Bekki}
     {Analysis and Experimental Verification of Weak Crossover in the Double Object and Oblique Dative Constructions}
     {pp.73--78}
  
\talk{Lisa Bylinina}
     {Augustinian Language Learning: What Ostensive Definition Can and Cannot Teach a Small Language Model}
     {pp.79--84}



\slot{15:30-17:00}{Expressives and evaluativity}

  
\talk{Linmin Zhang}
     {What-exclamatives, indescribability, and extremeness}
     {pp.85--90}
  
\talk{Osamu Sawada}
     {Two types of rhetorical wh-expressives in Japanese}
     {pp.91--96}
  
\talk{Frank Sode}
     {The Use Conditions of Melioratives and Tense/Mood-Marking}
     {pp.97--100}



\slot{17:30-18:30}{Quantification and genericity}

  
\talk{Emil Eva Rosina, Michaela Hulmanová and Marcin Wągiel}
     {Events are times}
     {pp.101--104}
  
\talk{Satoru Suzuki}
     {Intuitionistic-Bayesian Semantics of Generics Based on Belief Contraction}
     {pp.105--110}



\slot{19:30-}{Conference dinner}




\noindent\textbf{\large 
3rd Day: November 28 (Sat)
}\\



\slot{10:00-11:30}{Vagueness and degrees}

  
\talk{Ka-Fat Chow}
     {A Boolean Many-Valued Treatment of Vague Quantified Statements}
     {pp.111--116}
  
\talk{Linmin Zhang and Hanxi Zheng}
     {Estimating the positive threshold of gradable adjectives with implicit comparison}
     {pp.117--122}
  
\talk{Yi-Hsun Chen}
     {A degree-semantic analysis of possessive measure phrases}
     {pp.123--126}



\slot{13:30-15:00}{Atoms, Units, and Counting}

  
\talk{Yi-Hsun Chen and Ryan Smith}
     {Two types of nominal similatives: the case of Mandarin Chinese}
     {pp.127--130}
  
\talk{Ruoying Zhao, Abdullah Almuhaysh, Victoria Noble, Luca Molinari and Marcin Wągiel}
     {What a mass! Counting Hasawi Arabic mass nouns without coercion}
     {pp.131--135}
  
\talk{Shunsuke Yatabe}
     {Trinity}
     {pp.136--140}



\slot{15:30-17:00}{Incrementality and scalar particles}

  
\talk{Yu-Lun Huang and I-Ta Chris Hsieh}
     {On the incremental use of again: a view from Mandarin}
     {pp.141--145}
  
\talk{Chungmin Lee and Dongsik Lim}
     {Incrementality and even-like particles in Korean}
     {pp.146--149}
  
\talk{Ruoying Zhao and Moreno Mitrović}
     {Almost at-issue negative inference licenses expletive negation and more}
     {pp.150--155}



\slot{17:30-18:30}{Invited Talk}

  
\talk{Daniel Lassiter}
     {Presupposition, domain restriction and trivalent conditionals}
     {p.156}



\slot{18:30-18:40}{Closing remark}




  \end{document}

